
% tip to kill the wrong type
%!TEX program = xelatex
\documentclass[10pt,journal,compsoc]{IEEEtran}
\IEEEoverridecommandlockouts
% The preceding line is only needed to identify funding in the first footnote. If that is unneeded, please comment it out.
\usepackage{cite}
\usepackage{amsmath,amssymb,amsfonts}
\usepackage{algorithmic}
\usepackage{graphicx}
\usepackage{textcomp}
\usepackage{xcolor}
%\usepackage{}
\usepackage{algorithm}
\usepackage{xeCJK}
\def\BibTeX{{\rm B\kern-.05em{\sc i\kern-.025em b}\kern-.08em
    T\kern-.1667em\lower.7ex\hbox{E}\kern-.125emX}}
\usepackage{fancyhdr}
\begin{document}
kErrb3rt sCh1ld

\section{第一章}
*\textcolor{red}{基础章节}

1:坐标表象:
\begin{align}
  \hat p=-i\hbar \nabla\quad \hat H =-\frac{\hbar^2}{2\mu}\nabla^2+V
\end{align}
以及薛定谔方程:
\begin{align}
  \hat H\phi=i\hbar \frac{ \partial }{ \partial t}\psi
\end{align}

2:傅里叶变换
\begin{align}
  & \psi(p)=\frac1{\sqrt{2\pi\hbar}}\int \psi(x)e^{-ipx/\hbar}dx \\
  &\psi(x)=\frac{1}{\sqrt{2\pi\hbar}}\int \psi(p)e^{ipx/\hbar}dp
\end{align}

3:传播子问题
\begin{align}
  \psi(t)=\Sigma_nC_n\psi_n(0)e^{-iEt/\hbar}
\end{align}
其中，$C_n$的求法:\begin{align}
  C_n=\psi^*_n\cdot\psi(0)
\end{align}

4:定态和非定态

(a):概率密度和概率流密度不随时间改变。

(b):不含t的力学量F的平均值不随时间改变。

(c):力学概率分布随时间不变。

(d):\textcolor{red}{束缚态为非简并能量本征态。}

5:不确定关系
\begin{align}
  \Delta u=\sqrt{\overline{u^2}-{\overline{u}}^2}
\end{align}

6:积分公式\textcolor{red}{多次使用过！！！}
\begin{align}
&\int^{\infty}_{0} x^me^{{-\alpha x}}dx = \frac{m!}{\alpha^{m+1}} \\
& \int^{\infty}_{-\infty}x^{2n}e^{-\alpha x^2}dx=\frac{1\cdot3\cdot5\dots (2n-1)}{(2\alpha)^n}\sqrt{\frac{\pi}{\alpha}}\\
&\int^{\infty}_{-\infty}e^{-\alpha x^2+bx}dx=\sqrt{\frac{\pi}{\alpha}}e^{\frac{b^2}{4\alpha}}
\end{align}

\section{第二章}
*\textcolor{red}{重要章节，除了简谐要使用9章的算子算}

1:在V确定后，解Shr\"odinger方程


2:宇称相关定理：若V(x)=V(-x)，则有宇称确定，则解出的波函数不可能是sin和cos的结合。只能是二者之一，最后就会造成$\psi(x)=\psi(-x)$或者$\psi(x)=-\psi(-x)$。

3:关于透射率的计算

若有\begin{align}
  \psi=\begin{cases}
& e^{ikx}+Re^{-ikx}\quad x<0\\
&Se^{ik'x}\quad x>0
\end{cases}
\end{align}
最后求得的反射率等于$|R|^2$,这个是大概率成立的，因为反射后的k和反射前的k是相等的，但是透射率T不一定等于$|S|^2$，因为投射后的k'和透射之前不同，而T和$|S|^2$的关系有：\begin{align}T=|S|^2\frac{k'}k\end{align}，所以\textcolor{red}{直接求反射律，然后用1-反射律求出透射律。}而折射律和投射律的准确定义为:
\begin{align}
  & T=\frac{j_t}{j_i}\\
  & \text{反射律}=\frac{j_r}{j_i}\\
  &j=\frac1{2\mu}[\psi^*\hat p\psi-\psi\hat p\psi^*]
\end{align}
在这个公式代入$\psi_i=e^{ikx}$，$\psi_t=Te^{ik'x}$，和$\psi_r=Re^{-ikx}$。就得到了$j_i,j_t,j_r$。

4:关于简谐的计算：用第九章的算子和粒子数态计算。

\section{第三章}
*\textcolor{red}{过渡章节}

1:角动量和动量相关的关系式
\begin{align}
  & [l_i,l_j]=i\hbar \epsilon_{ijk}l_k \\
  & [l_i,p_j]=i\hbar \epsilon_{ijk}p_k\\
  & [l_i,x^j]=i\hbar x^k\epsilon_{ijk}\\
  & [x,p]=i\hbar\\
\end{align}

2:升降算符
\begin{align}
  j_{\pm}=l_x\pm il_y
\end{align}


3:久期方程的使用

久期方程就是能量本征方程在表象下的分量表达式。

4:泊松括号和对易的关系

\begin{align}
  [A,B]=i\hbar\{\frac{ \partial A}{ \partial q}\frac{ \partial B}{ \partial p}-\frac{ \partial A}{ \partial p}\frac{ \partial B}{ \partial q} \}
\end{align}
注意这里的A和B只能是单项，不能是$x^2$。

\section{第四章}
*\textcolor{red}{过渡章节}

1:多粒子体系的波函数的特征

(a):Fermion:半自旋，多粒子体系时，两个粒子的位置交换，波函数反对称。

(b):Boseon:整自旋，多粒子体系时，两个粒子的位置交换，比函数对称。

(c):假设有两个粒子：存在可能的两个态:$\psi_1,\psi_2$。并且默认写在前表示第一个粒子，写在后表示第二个粒子，则$\psi_1(1)\psi_2(2)=\psi_{1}\psi_2,\psi_1(2)\psi_2(1)=\psi_{2}\psi_1$

(i):若两个粒子为Boseon,则成立的状态有$\psi_{11},\psi_{22},\psi_{12}$,其中$\psi_{11},\psi_{22}$本身就是对称的，但是$\psi_{12}$，不是对称的，则实际存在的状态必须要对称化\begin{align}
  \psi_{12}=\frac1{\sqrt{2}}[\psi_{1}\psi_2+\psi_2\psi_1]
\end{align}

(ii):同样的，对于Fermion,像$\psi_{12}$的状态要反对称处理，并且$\psi_{11},\psi_{22}$不存在。

2:\textcolor{red}{海森堡方程}

\begin{align}
  \frac{dA}{dt}=\frac1{i\hbar}[A,H]
\end{align}

3:守恒量的推导

\begin{align*}
& \frac{d\overline{A}}{dt}=\frac{d}{dt}\int \psi^* A\psi=\int \frac{d}{dt}\psi^* A \psi+\int \psi^* \frac{dA}{dt}\psi+\int \psi^* A \frac{d\psi}{dt}\\
& = \frac1{-i\hbar}[\int H\psi^* A\psi-i\hbar\int\psi^*\frac{d}{dt} A\psi-\int \psi^* A H\psi ]\\
&=\frac1{i\hbar}\overline{[A,H]}+ \int{\psi^*\frac{d}{dt} A\psi}
\end{align*}
若A不随时间变化，则得到守恒量的表达。

4:位力定理\begin{align}
  2\overline{T}=\overline{r\cdot\nabla V}
\end{align}

\section{第七章}
*\textcolor{red}{过渡章节}

1:本征态对应的矩阵为对角矩阵。

2:注意矩阵运算(矩阵证明题类似作业唯二题目)。

\section{第八章}
*\textcolor{red}{重要章节！！！}

1:电子自旋算符和Pauli算符的关系:
\begin{align}
  s=\frac12\hbar\sigma
\end{align}

2:在z方向自旋表象下的Pauli算符:
\begin{align}
  \sigma_x=\begin{bmatrix}0&1\\1&0\end{bmatrix}\quad
  \sigma_y=\begin{bmatrix}0&-i\\i&0\end{bmatrix}\\
  \sigma_z=\begin{bmatrix}1&0\\0&-1\end{bmatrix}
\end{align}注意这个关系知道的比$\sigma_{\pm}$更早且更简单。由\textcolor{red}{简答推复杂关系也可以记忆}
定义$\sigma_{\pm}=\frac12(\sigma_x\pm i \sigma_y)$。

3:双粒子问题:
(a):若双Boseon,则波函数必定为对称，若某个态为$\psi_{12}$则要进行对称化，把波函数化为：\begin{align}
  \psi_{12}=\frac1{\sqrt{2}}(\psi_1\psi_2+\psi_2\psi_1)
\end{align}默认写在前为第一个粒子。而$\psi_{11}$就不需要对称化。

(b):若为双Fermion，则波函数必定为反称的，不存在态$\psi_{nn}$,对于态$\psi_{12}$要进行反称化：

\begin{align}
  \psi_{12}=\frac1{\sqrt{2}}(\psi_1\psi_2-\psi_2\psi_1 )
\end{align}

4:双Fermion自旋耦合的自旋算符:\begin{align}\hat S=\hat s_1+\hat s_2\end{align}对于这个算符，可以选\textcolor{red}{**三**}种基底:

(a):$(S^2,S_z) $本征态基底，为\begin{align}
&  \chi_{00}=\frac1{\sqrt{2}}(\chi_1\chi_0-\chi_0\chi_1)\\
&  \chi_{1-1}=\chi_0\chi0\\
&  \chi_{10}=\frac1{\sqrt{2}}(\chi_1\chi_0+\chi_0\chi_1)\\
&  \chi_{11}=\chi_1\chi_1
\end{align}
通式表达为:$\chi_{Sm_s}$，并且也满足角动量的表达式。

(b):非本征态基底\begin{align}
\chi_1\chi_1,\chi_1\chi_0,\chi_0\chi_1,\chi_0\chi_0
\end{align}

(c):Bell基:\begin{align}
  (\sigma_{1z}\sigma_{2z},\sigma_{1x}\sigma_{2x})
\end{align}的共同本征基底：

\begin{align}
  & \psi^+_{12}=\frac1{\sqrt{2}}(\chi_1\chi_0+\chi_0\chi_1)\\
  &\psi^-_{12}=\frac1{\sqrt{2}}(\chi_1\chi_0-\chi_0\chi_1) \\
  &\phi^+_{12}=\frac1{\sqrt{2}}(\chi_1\chi_1+\chi_0\chi_0) \\
  &\phi^-_{12}=\frac1{\sqrt{2}}(\chi_1\chi_1-\chi_0\chi_0)
\end{align}简单记忆：$\psi$和$\phi$用于区分同向和反向，正负号用于区分对称和反称。

***\textcolor{red}{不论选何种基底，讨论问题时首先关注的是对称和反称。}***

\section{第五章}
*\textcolor{red}{几乎不考,重点看1,b就行}

1:有心势能的波函数:

(a):解有心势能的Shr\"odinger方程:
\begin{align}
  &[-\frac{\hbar^2}{2\mu r^2}\frac{ \partial }{ \partial r}r^2\frac{ \partial }{ \partial r}+\frac{l(l+1)\hbar^2}{2\mu r^2}+V(r)]R(r)=ER(r)\\
  & \chi(r)=R\cdot r\\
  &[\frac{d^2\chi}{dr^2}+[\frac{2\mu}{\hbar^2}(E-V(r))-\frac{l(l+1)}{r^2}]\chi(r)]=0
\end{align}

(b)\textcolor{red}{氢原子}

对于氢原子的势能$V=-\frac{e^2}{r}$，在把这个表达式代入39式后，在极限的情况下解出的波函数为:
\begin{align}
  & r\rightarrow 0 :\psi \rightarrow r^l\\
  & r \rightarrow \infty :\psi \rightarrow e^{-\beta r}\\
  &\text{对于一般的r} \psi = r^le^{-\beta r}u(r)\\
  &\text{ 其中的参数:}\beta=\frac1n=\sqrt{-2E}(\text{在自然单位下})\\
\end{align}
同时，在条件$\alpha=-n_r=l+1-n=0$(同时也是轨道为圆轨道的条件)的条件下，u(r)=1。

2:参数意义:

$\alpha=-n_r$其中$n_r$为径向量子数，衡量镜像的节点。

3:\textcolor{red}{$E_n$的简并度是$n^2$。}长度单位:$L=\frac{\hbar^2}{\mu e^2} $。能量单位为:$E=\frac{\mu e^4}{\hbar^2} $

\section{第九章}
*\textcolor{red}{重要章节}

1:粒子数和产生和湮灭算符，以及粒子数算符:
\begin{align}
  & a = \frac1{\sqrt{2}}(x+ip),a^{\dagger}=\frac1{\sqrt{2}}(x-ip)\\
  & \hat N=a^{\dagger}a
\end{align}
对于本征波函数$|n>$作用后有:
\begin{align}
  & a^{\dagger}|n>=\sqrt{n+1}|n+1>,a|n>=\sqrt{n}|n-1>\\
  & \hat N |n>=n|n>
\end{align}

2:对易关系式:
\begin{align}
  & [a,a^{\dagger}]=1,[a,a^{\dagger}a]=a
\end{align}

3:多粒子的代数角动量:双简谐振子定义:
\begin{align}
    & \hat j_+=a^{\dagger}_1a_2,\hat j_-=a_1a^{\dagger}_2\\
    & \hat j_z=\frac12(\hat N_1-\hat N_2)\\
    & j=\frac12n=\frac12(n_1+n_2) , m=\frac12(n_1-n_2)
\end{align}

\section{第十章}
*\textcolor{red}{中等}

1:非简并微扰论:计算步骤

(a):把H写为$H^{(0)}+H'$。

(b):求$\psi^{(0)}_n$，并且求出$H'_{nn}$。知道$E^{(1)}_n=H'_{nn}$。

(c):再求$E^{(2)}_n=\Sigma_m\frac{|H'_{mn}|^2}{E^{(0)}_n-E^{(0)}_m}$。并且得到:$\psi^{(1)}_n=\Sigma_m\frac{|H'_{mn}|}{E^{(0)}_n-E^{(0)}_m}\psi^{(0)}_m $ 。

(d):最后得到总和的:
\begin{align}
  & E_n=E^{(0)}_n+H'_{nn}+\Sigma_m\frac{|H'_{mn}|^2}{E^{(0)}_n-E^{(0)}_m}\\
  & \psi_n=\psi_n^{(0)}+\Sigma_m\frac{|H'_{mn}|}{E^{(0)}_n-E^{(0)}_m}\psi_m^{(0)}
\end{align}

2:简并态的微扰论:

(a):把H写为$H^{(0)}+H'$。并且解出$\psi_n^{(0)}$。

(b):解关于$H'$的久期方程，解出$\lambda'_n=E^{(1)}_n$,以及$\psi'_n$其中，$\lambda'_n$为能量的修正值，并且$\psi'_n$为新的波函数。


3:近简并态的微扰论:

(a):原先的Hamiltonian为$\begin{bmatrix}H_{11}&0\\0&H_{22}\end{bmatrix} $，并且$H_{11}-H_{22}\sim 0$，这这种情况下，对于任意微扰:$\begin{bmatrix}0&H'_{12}\\H'_{21}&0\end{bmatrix}  $先得到完整的Hamiltonian。

(b):直接用久期方程求完整的Hamiltonian的本征值本征态。

(c):讨论H'对结果的影响。
\section{一些技巧}
这些技巧是从经验里得来的，可以用于偷懒。

1:易知推难知原则：

在Pauli和自旋的关系中，我们直到的是$\sigma_x$而不是$\sigma_{\pm}$,相对应的，在角动量的关系中，我们直到的是$l_{\pm}$，在简谐粒子数表象中，易知的是$a$和$a^{\dagger}$，总结出易求推难求原则：
\begin{align}
&
  \sigma_{\pm}=\frac12(\sigma_x\pm i\sigma_y)
\\
&\begin{cases}
  &l_x=\frac12(l_++l_-)\\
  &l_y=\frac1{2i}(l_+-l_-)
\end{cases}\\
&\begin{cases}
  & x=\frac1{\sqrt{2}}(a^{\dagger}+a)\\
  & p=\frac1{\sqrt{2}i}(a-a^{\dagger})
\end{cases}
\end{align}

2:通用公式:
\begin{align}
  j_{\pm}\psi=\sqrt{j(j+1)-m(m\pm1)}\hbar \psi\\
  j_z\psi=m\hbar\psi
\end{align}
33式为通用的表达式，34式为使用条件。下面是一些例子:
(a):\begin{align}
  l_+Y_{lm}=\sqrt{l(l+1)-m(m+1)}\hbar Y_{lm+1}
\end{align}

(b):\begin{align}
  j_-Y_{lm}=\sqrt{l(l+1)-m(m-1)}\hbar Y_{lm-1}
\end{align}

(c):\textcolor{red}{有点重要}:\begin{align}
  s_+\chi_{-\frac12}=\sqrt{\frac34+\frac14}\hbar \chi_{\frac12}
\end{align}
c的解释：把$\chi_{\pm\frac12}$看作$Y_{\frac12\pm\frac12}$代入通用公式。

\newpage

\section{所有用到过的公式}

\subsection{第一章}

\begin{align*}
  &-\frac{\hbar^2}{2m}\nabla^2\psi+V\psi=i\hbar\frac{ \partial }{ \partial t}\psi=E\psi \\
  & \int pdq=nh \\
  &\textcolor{red}{\Delta u=\sqrt{\overline{u^2}-\overline{u}^2}} \\
  &\psi(p)=\frac1{\sqrt{2\pi\hbar}}\int\psi(x)e^{-ipx/\hbar}dx,\\ &\psi(x)=\frac1{\sqrt{2\pi\hbar}}\int \psi(p)e^{ipx/\hbar}dp \\
  &\textcolor{red}{ C_n=\int \psi^*_n \psi \quad ;\;\psi(t)=\Sigma_n C_n\psi_n e^{-iEt/\hbar}} \\
  &\textcolor{red}{ \int^{\infty}_{-\infty}e^{-\alpha x^2}dx=\sqrt{\frac{\pi}{\alpha}} }\\
  &\textcolor{red}{ \int^{\infty}_{-\infty}x^2e^{-\alpha x^2}dx=\frac{1}{2\alpha}\sqrt{\frac{\pi}{\alpha}}} \\
\end{align*}

\subsection{第二章}

\begin{align*}
  & \psi=\begin{cases}e^{ikx}+Re^{-ikx}\\ Se^{ik'x}\end{cases};T\ne |S|^2,T+|R|^2=1 \\
\end{align*}

\subsection{第三章}

\begin{align*}
  &\textcolor{red}{ [l_i,A_j]=i\hbar\epsilon_{ijk}A_k,A=l,x,p} \\
  &\textcolor{red}{ j_{\pm}=j_x\pm ij_y }\\
  &\textcolor{red}{ j_{\pm}|nlm>=\sqrt{l(l+1)-m(m\pm1)}\hbar|nlm\pm1> } \\
  &\textcolor{red}{ \text{久期方程}} \\
  &[A,B]=i\hbar(\frac{ \partial A}{ \partial q}\frac{ \partial B}{ \partial p}-\frac{ \partial A}{ \partial p}\frac{ \partial B}{ \partial q}) \\
\end{align*}

\subsection{第四章}

\begin{align*}
  &\textcolor{red}{ \frac{dA}{dt}=\frac1{i\hbar}[A,H]=\frac1{-\hbar^2}\{A,B\} }\\
  &2\overline{T}=\overline{r\cdot \nabla V} \\
\end{align*}

\subsection{第七章}

没有关键公式

\subsection{第八章}

\begin{align*}
  &\sigma_x=\begin{bmatrix}0&1\\1&0\end{bmatrix},\sigma_y=\begin{bmatrix}0&-i\\i&0\end{bmatrix}  \\
  &\sigma_z=\begin{bmatrix}1&0\\0&-1\end{bmatrix} ,\textcolor{red}{ \sigma_\pm=\frac12(\sigma_x\pm i\sigma_y)} \\
\end{align*}

\subsection{第五章}

\begin{align*}
&\text{$n_l=0$时，有}u(r)=1\rightarrow R_{nl}=r^le^{-\beta r} \\
&\chi(r)=r^{l+1}e^{-\beta r},\beta =\frac1n,n=l+1 \\
\end{align*}


\subsection{第九章}

\begin{align*}
    & \textcolor{red}{ a|n>=\sqrt{n}|n-1>,a^{\dagger}|n>=\sqrt{n+1}|n+1> }\\
    &\hat N=a^{\dagger}a,aa^{\dagger}=\hat N+1 \\
    &\textcolor{red}{ x=\frac1{\sqrt{2}}(a+a^{\dagger}),p=\frac1{\sqrt{2}i}(a-a^{\dagger}) }\\
    &\text{x和p的单位是：}\textcolor{red}{ |x|=\sqrt{\frac{\hbar}{\mu\omega}},|p|=\sqrt{\hbar\omega\mu} }\\
  \end{align*}

  \subsection{第十章}

  非简并微扰:
  \begin{align*}
    &E_n=E^{(0)}_n+H'_{nn}+\textcolor{red}{ \Sigma_m\frac{|H'_{mn}|^2}{E^{(0)}_n-E^{(0)}_m}}\\
    &\psi_n=\psi^{(0)}_n+\Sigma_m\frac{|H'_{mn}|}{E^{(0)}_n-E^{(0)}_m}\psi^{(0)}_m
  \end{align*}

  简并态微扰:

  \begin{align*}
  & det|H'-\lambda|=0\rightarrow v_n,\lambda_n\\
  &\psi'_n=v_n,E_n=E_n^{(0)}+\lambda_n\\
  \end{align*}

  逼近简并的非简并:

  久期方程

\end{document}
